% Options for packages loaded elsewhere
\PassOptionsToPackage{unicode}{hyperref}
\PassOptionsToPackage{hyphens}{url}
\PassOptionsToPackage{dvipsnames,svgnames,x11names}{xcolor}
%
\documentclass[
  11pt,
  a4paper,
]{ltjsarticle}
\usepackage{amsmath,amssymb}
\usepackage{iftex}
\ifPDFTeX
  \usepackage[T1]{fontenc}
  \usepackage[utf8]{inputenc}
  \usepackage{textcomp} % provide euro and other symbols
\else % if luatex or xetex
  \usepackage{unicode-math} % this also loads fontspec
  \defaultfontfeatures{Scale=MatchLowercase}
  \defaultfontfeatures[\rmfamily]{Ligatures=TeX,Scale=1}
\fi
\usepackage{lmodern}
\ifPDFTeX\else
  % xetex/luatex font selection
\fi
% Use upquote if available, for straight quotes in verbatim environments
\IfFileExists{upquote.sty}{\usepackage{upquote}}{}
\IfFileExists{microtype.sty}{% use microtype if available
  \usepackage[]{microtype}
  \UseMicrotypeSet[protrusion]{basicmath} % disable protrusion for tt fonts
}{}
\makeatletter
\@ifundefined{KOMAClassName}{% if non-KOMA class
  \IfFileExists{parskip.sty}{%
    \usepackage{parskip}
  }{% else
    \setlength{\parindent}{0pt}
    \setlength{\parskip}{6pt plus 2pt minus 1pt}}
}{% if KOMA class
  \KOMAoptions{parskip=half}}
\makeatother
\usepackage{xcolor}
\usepackage[margin=24mm]{geometry}
\setlength{\emergencystretch}{3em} % prevent overfull lines
\providecommand{\tightlist}{%
  \setlength{\itemsep}{0pt}\setlength{\parskip}{0pt}}
\setcounter{secnumdepth}{5}
\ifLuaTeX
  \usepackage{selnolig}  % disable illegal ligatures
\fi
\IfFileExists{bookmark.sty}{\usepackage{bookmark}}{\usepackage{hyperref}}
\IfFileExists{xurl.sty}{\usepackage{xurl}}{} % add URL line breaks if available
\urlstyle{same}
\hypersetup{
  colorlinks=true,
  linkcolor={blue},
  filecolor={Maroon},
  citecolor={Blue},
  urlcolor={blue},
  pdfcreator={LaTeX via pandoc}}

\author{}
\date{}

\begin{document}

\hypertarget{ux6700ux5c0fux4e8cux4e57ux6cd5}{%
\section{最小二乗法}\label{ux6700ux5c0fux4e8cux4e57ux6cd5}}

方程式

\[
Ax=b
\]

が厳密に解けない場合、最小二乗法では

\[
\boxed{\min_x\|Ax-b\|_2^2}
\]

を解きます。

\hypertarget{ux5e7eux4f55ux5b66}{%
\subsection{幾何学}\label{ux5e7eux4f55ux5b66}}

\texttt{Ax} は必ず \texttt{Col(A)}
上にあります。したがって問題は、\texttt{b}
に最も近い列空間上の点を探すことです。

最適点を \texttt{A\textbackslash{}hat\ x} とすると、残差

\[
r=b-A\hat x
\]

は列空間に直交します。

\[
A^\top r=0
\]

よって

\[
A^\top(b-A\hat x)=0
\]

から

\[
A^\top A\hat x=A^\top b
\]

が得られます。

\hypertarget{ux7d71ux8a08ux3068ux306eux63a5ux7d9a}{%
\subsection{統計との接続}\label{ux7d71ux8a08ux3068ux306eux63a5ux7d9a}}

線形回帰

\[
y=X\beta+\varepsilon
\]

で二乗誤差を最小化することは、\texttt{y} を \texttt{X}
の列空間へ射影することと同じです。

つまり最小二乗法は「回帰の公式」ではなく、直交射影という線形代数の幾何学です。

\end{document}
